\vspace{-0.4cm}
\section{Introduction}
Large Language Models (LLMs) have demonstrated impressive reasoning abilities, such as in solving complex math problems. A key factor driving this progress is the use of Chain-of-Thought (CoT) reasoning \citep{wei2022chain}, where the model engages in intermediate deliberation before producing an answer. Building on this, recent advances have employed Reinforcement Learning (RL) to further enhance LLM reasoning by optimizing for verifiable outcome rewards \citep{jaech2024openai, guo2025deepseek, yu2025dapo, wei2025swe}. Notably, RL-trained models have exhibited emergent behaviors such as engaging in self-reflection, a process of revisiting previous states to correct earlier mistakes \citep{guo2025deepseek, zeng2025simplerl}, which offers a potential way for test-time scaling \citep{snell2024scaling,brown2024large} by generating longer CoTs. However, despite these compelling phenomena, it remains unclear why and under what conditions self-reflections are beneficial, or whether such behaviors will emerge through conventional RL training.

In this work, we use reflective exploration to refer to behavior where the model revisits a previously visited state and, informed by the additional context it has acquired, alters its subsequent actions. This definition links self-reflective directly to exploration under uncertainty, providing a principled basis for when and why strategy switches should occur rather than treating them as ad hoc quirks of the model’s CoTs. We formalize this via strictly uncertainty-adaptive policies whose actions depend nontrivially on the belief. Since conventional RL typically admits a Markovian optimal policy, reflective exploration is not induced since the policy depends on the history only through the state and is thus not incentivized to enrich identical states with additional context.
Instead, RL explores during training through repeated trial and error, with no incentives for further exploration when parameters are frozen after deployment. Therefore, under conventional RL, there is no guarantee that self-reflections will emerge during training, nor does it explain why they might be advantageous when the trained policy is deployed at test time.

To address this gap, we ground reflective exploration with Bayesian RL, which maximizes the expected return in the deployment stage under a posterior distribution over Markov Decision Processes (MDPs) induced by the training data. The objective incentivizes both reward-seeking actions and exploration actions that gather information to reduce the MDP's uncertainty, such as the reward uncertainty regarding the progress made by different actions. 
We show that optimizing this Bayesian RL objective yields policies that adapt on-the-fly by updating their beliefs and switching strategies based on observed outcomes, naturally leading to reflective exploration behaviors. Moreover, we prove that the optimal adaptive policy can achieve arbitrarily higher Bayes-expected return than any optimal Markovian policy.

Building upon this formulation, we introduce a novel algorithm, \textit{Bayes-Adaptive RL for LLM Reasoning} (BARL). For each prompt, BARL performs online rollouts to generate a set of candidate answers, each associated with an MDP hypothesis. 
The state-action value is then computed by weighting each hypothesis according to the model’s current belief, with penalties applied for mismatches between predicted and observed rewards, thereby signaling when to switch strategies. BARL provides a principled mechanism for integrating and revising plausible strategies, analogous to linearizing best-of-N, but with explicit guidance on \textit{when} and \textit{how} the model should reflectively explore.

To illustrate the benefits of BARL, we begin with a synthetic task designed to mirror the training-deployment procedure in LLM reasoning. The agent receives a reward only when it repeats a prompt token exactly three times, but the training and evaluation prompt tokens differ. Conventional RL fails to generalize beyond the training prompt tokens. In contrast, BARL learns to switch strategies by eliminating hypotheses, ultimately discovering the ground-truth MDP for optimal behavior.

We further evaluate BARL on math reasoning tasks using various LLMs, including Qwen2.5-Math-1.5B, Qwen2.5-Math-7B, and R1-Distill-Llama-8B. Across these models, BARL consistently outperforms conventional RL algorithms, such as GRPO and a strong progress-reward baseline, on multiple benchmarks. BARL achieves significantly greater token efficiency, requiring up to \textbf{1.63x} fewer average tokens than the progress baseline, \textbf{2x} fewer than GRPO, and over \textbf{10x} fewer than the Qwen2.5-Math-1.5B base model. Moreover, we observe no strong correlation between overall performance and the frequency of reflections. Instead, BARL’s advantage stems from more efficient exploration and more effective thinking tokens.

\vspace{0.15cm}
\begin{AIbox}{\text{Key Takeaways: Why, How, and When Should LLMs Reflectively Explore}}
\begin{itemize}[leftmargin=0em]
\item \textbf{Why:} RL neither ensures the emergence of reflective exploration nor explains its benefits since (1) RL attains Markovian optimal policies, which cannot exhibit reflective exploration that collects additional context and act differently at the same state, and (2) RL exploration is confined to trial-and-error during training, leaving no mechanism that encourages adaptive exploration at test time. In contrast, Bayesian RL, by optimizing Bayes-expected return during deployment, encourages exploration to gather contextual information that reduces the MDP uncertainty.
\item \textbf{How:} BARL provides a principled way to stitch plausible strategies by maintaining a posterior over MDP hypotheses, each associated with a sampled candidate answer. Reflective exploration emerges through hypothesis elimination, enabling on-the-fly adaptation.
\item \textbf{When:} LLMs should self-reflect when discrepancies arise between their internal beliefs and cumulative reward feedback—signaling strategy switching by downweighting hypotheses that have high belief probabilities but are unlikely to be optimal given previous observations.
\end{itemize}
\end{AIbox}
